\section*{Chapitre \ref{ch:g12:randvar} --- Variables aléatoires et loi binomiale}

\begin{solution}{exo:g12:randvar:1}
$\E(X) = -0.3 + 0 + 0.8 + 0.5 = 1$.
$\E(X^2) = 0.3 \times 1 + 0 + 0.4\times4 + 0.1\times25 = 4.4$, donc par
König--Huygens $\V(X) = 4.4 - 1 = 3.4$ et
$\sigma(X) = \sqrt{3.4} \approx 1.84$.
\end{solution}

\begin{solution}{exo:g12:randvar:2}
\emph{1.} Les $10$ questions sont des \omterm{def:g12:randvar:bernoulli}{épreuves de Bernoulli} indépendantes
avec $p = \frac14$~: $X \sim \mathcal B\left(10, \frac14\right)$,
$\E(X) = 2.5$,
$\sigma(X) = \sqrt{10 \times \frac14 \times \frac34} = \sqrt{1.875}
\approx 1.37$.

\emph{2.} $\P(X = 0) = \left(\frac34\right)^{10} \approx 0.056$~;
\[
\P(X = 5) = \binom{10}{5}\left(\frac14\right)^5\left(\frac34\right)^5
\approx 0.058;
\quad
\P(X \geq 1) = 1 - \left(\tfrac34\right)^{10} \approx 0.944 .
\]
\end{solution}

\begin{solution}{exo:g12:randvar:3}
Essais identiques \omterm{def:g12:condprob:indep}{indépendants}~: $X \sim \mathcal B(400,\ 0.05)$, donc
$\E(X) = 20$ réclamations et
\[
\sigma(X) = \sqrt{400 \times 0.05 \times 0.95} = \sqrt{19} \approx 4.4 .
\]
\end{solution}

\begin{solution}{exo:g12:randvar:4}
Le paiement $P$ reçu vérifie
$\E(P) = \frac{5 + 6}{6} = \frac{11}{6}$, donc le prix équitable est
$c = \frac{11}{6} \approx 1.83$~euros. Le gain équitable
$G = P - \frac{11}6$ prend les valeurs $-\frac{11}{6}, \frac{19}{6},
\frac{25}{6}$ avec \omterm{def:g10:proba:distribution}{probabilités} $\frac46, \frac16, \frac16$~:
\[
\V(G) = \E(G^2)
= \frac{4 \times 121 + 361 + 625}{6 \times 36} = \frac{1470}{216}
= \frac{245}{36} \approx 6.8,
\qquad \sigma(G) \approx 2.6 .
\]
Un jeu équitable a un gain \emph{moyen} nul mais ses issues fluctuent
encore~: équitable n'est pas sans risque.
\end{solution}

\begin{solution}{exo:g12:randvar:5}
$X \sim \mathcal B(8,\ 0.7)$.
\[
\P(X = 6) = \binom86 (0.7)^6(0.3)^2 \approx 0.296 ;
\]
\[
\P(X \geq 6) = \P(6) + \P(7) + \P(8)
\approx 0.296 + 8(0.7)^7(0.3) + (0.7)^8
\approx 0.296 + 0.198 + 0.058 = 0.552 ;
\]
$\P(X \geq 1) = 1 - (0.3)^8 \approx 0.99993$. Mode~:
$(n+1)p = 6.3$, donc la valeur la plus probable est $k^* = 6$
(\cref{exo:g12:randvar:7}).
\end{solution}

\begin{solution}{exo:g12:randvar:6}
Le nombre de passagers se présentant est $X \sim \mathcal B(104,\ 0.9)$~;
le vol est en surréservation lorsque $X \geq 101$~:
\[
\P(X \geq 101) = \sum_{k=101}^{104}\binom{104}{k}(0.9)^k(0.1)^{104-k}
\approx 0.006 .
\]
Vendre $4\%$ de billets de plus que de places provoque un incident sur
seulement environ $0.6\%$ des vols --- l'économie derrière la
surréservation.
\end{solution}

\begin{solution}{exo:g12:randvar:7}
\[
\frac{\P(X = k+1)}{\P(X = k)}
= \frac{\binom{n}{k+1}}{\binom nk}\cdot\frac{p}{1-p}
= \frac{n - k}{k + 1}\cdot\frac{p}{1-p},
\]
en utilisant $\binom{n}{k+1} = \binom nk \frac{n-k}{k+1}$. Ce rapport est
$\geq 1$ ssi $(n-k)p \geq (k+1)(1-p)$ ssi $np - k p \geq k - kp + 1 - p$
ssi $k \leq (n+1)p - 1$. Donc les \omterm{def:g10:proba:distribution}{probabilités} croissent strictement tant
que $k + 1 \leq (n+1)p$ et décroissent ensuite~: le \omterm{def:g10:functions:extrema}{maximum} est atteint en
$k^* = \floor{(n+1)p}$ (partagé avec $k^* - 1$ lorsque $(n+1)p$ est un
entier).
\end{solution}

\begin{solution}{exo:g12:randvar:8}
\emph{1.} Pile d'abord au lancer $k$ signifie $k-1$ faces puis pile~:
\omterm{def:g10:proba:distribution}{probabilité} $\left(\frac12\right)^{k-1}\cdot\frac12 = 2^{-k}$, pour
$1 \leq k \leq 10$. \omterm{def:g10:proba:distribution}{Probabilité} totale de gagner
$\sum_{k=1}^{10} 2^{-k} = 1 - 2^{-10} = \frac{1023}{1024}$ (le jeu n'est
perdu que sur dix faces consécutives).

\emph{2.} Gain espéré~:
\[
\sum_{k=1}^{10} 2^k \cdot 2^{-k} = \sum_{k=1}^{10} 1 = 10 \text{ euros}.
\]
Sans le plafond, la somme $\sum_{k\geq1} 1$ diverge~: le gain espéré est
infini, bien que le jeu paie presque toujours un petit montant --- le
célèbre \emph{paradoxe de Saint-Pétersbourg}, montrant que l'\omterm{def:g12:randvar:exp}{espérance}
seule ne mesure pas la valeur d'un jeu.
\end{solution}

\begin{solution}{pb:g12:randvar:1}
\textbf{1.} $X \sim \mathcal B\!\left(3, \frac16\right)$~:
$\P(X = 0) = \frac{125}{216}$,
$\P(X = 1) = \frac{75}{216}$,
$\P(X = 2) = \frac{15}{216}$,
$\P(X = 3) = \frac{1}{216}$.

\textbf{2.} $\E(X) = 3 \times \frac16 = \frac12$~;
$V(X) = 3 \times \frac16 \times \frac56 = \frac{5}{12}$.

\textbf{3.} Versement espéré $2\,\E(X) = 1$ euro contre une mise
de $1$ euro~: gain nul --- parfaitement équitable (une rareté).

\textbf{4.} Sans remise, les tirages sont dépendants (les chances
de la seconde carte dépendent de la première)~: la \omterm{def:g12:randvar:rv}{loi} n'est pas
binomiale --- la case «~\omterm{def:g12:condprob:indep}{indépendance}~» de la liste de contrôle
échoue. Avec remise~: $n = 5$ fixé, $p = \frac14$ constant,
tirages \omterm{def:g12:condprob:indep}{indépendants}~: c'est binomial.

\textbf{5.} $\E(X) = 6$, $\sigma(X) = \sqrt{4.2} \approx 2.05$.
$\E(S) = 5 \times 6 - 10 = 20$~;
$\sigma(S) = 5\sigma(X) \approx 10.25$.

\textbf{6.} $n = 105$ épreuves fixées (les billets), chacune une
présence ou une absence de même paramètre $p = 0.9$, supposées
indépendantes~: c'est binomial. L'aveu~: les familles ratent leur
vol \emph{ensemble} et les tempêtes vident des avions entiers ---
l'\omterm{def:g12:condprob:indep}{indépendance} est l'acte de foi du modèle, et des absences
corrélées rendent les queues plus épaisses que ne le promet la
binomiale.

\textbf{7.} $\E(X) = 94.5$~;
$\sigma(X) = \sqrt{105 \times 0.9 \times 0.1} = \sqrt{9.45}
\approx 3.07$.

\textbf{8.} $\P(X \geq 101) = \sum_{k=101}^{105}
\binom{105}{k} 0.9^k\, 0.1^{105 - k}$.

\textbf{9.} Environ $0.0167$~: à peu près un vol sur soixante
refuse quelqu'un.

\textbf{10.} $\E(\text{refusés}) \approx 0.023$ passager par vol~:
indemnisation espérée d'environ $19$ euros --- contre $1\,000$
euros de recettes supplémentaires. Surréserver de cinq places est
massivement rentable~; d'où la pratique universelle.

\textbf{11.} D'après le tableau des \omterm{def:g10:proba:distribution}{probabilités} de queue~:
$n = 105$~: $1.7\,\%$~; $n = 106$~: $4.0\,\%$~; $n = 107$~:
$8.1\,\%$. La plus grande vente conforme est de $n = 106$ billets.

\textbf{12.} $z = \frac{100.5 - 94.5}{3.07} \approx 1.95$~: le
seuil se situe deux \omterm{def:g11:stat:variance}{écarts types} au-dessus de la \omterm{def:g11:stat:mean}{moyenne}, et la
règle empirique de la cloche («~queue supérieure au-delà de
$2\sigma$~$\approx 2.5\,\%$~») tombe près du $1.7\,\%$ exact ---
la courbe lisse qui double la binomiale est la protagoniste des
chapitres à venir.

\textbf{13.} $0.98^{20} + 20 \times 0.02 \times 0.98^{19}
\approx 0.94$~: le lot honnête passe $94\,\%$ du temps.

\textbf{14.} $0.9^{20} + 20 \times 0.1 \times 0.9^{19}
\approx 0.39$~: le mauvais lot se faufile encore $39\,\%$ du
temps.

\textbf{15.} Risque du producteur~: un bon lot refusé (environ
$6\,\%$)~; risque du consommateur~: un mauvais lot accepté (environ
$39\,\%$). Un \omterm{def:g10:proba:sample}{échantillon} plus grand (avec un seuil d'acceptation
proportionnel) réduit les deux --- au prix de davantage
d'inspection~: la qualité a sa ligne budgétaire.

\textbf{16.} $V\!\left(\frac Xn\right) = \frac{V(X)}{n^2} =
\frac{np(1-p)}{n^2} = \frac{p(1-p)}{n}$~: l'\omterm{def:g11:stat:variance}{écart type} vaut
$\sqrt{\frac{p(1-p)}{n}}$. Quadrupler l'\omterm{def:g10:proba:sample}{échantillon} divise le bruit
par deux~: la \omterm{def:g12:randvar:rv}{loi} en $\frac{1}{\sqrt n}$ des \omterm{def:g10:numbers:interval}{intervalles} de
fluctuation d'il y a deux ans, désormais démontrée.

\textbf{17.} Le plafond borne le gain à $2^{10} = 1024$ euros, si
bien que chacun des dix tours contribue exactement pour $1$ euro à
l'\omterm{def:g12:randvar:exp}{espérance}~: $\E = 10$. L'\omterm{def:g12:randvar:exp}{espérance} infinie du jeu non plafonné
venait entièrement de gains astronomiquement rares et
astronomiquement grands~; plafonner, c'est avouer qu'aucune banque
ne verse $2^{50}$ euros. Prix équitable du billet pour le jeu
plafonné~: $10$ euros --- et voilà le paradoxe dissipé.

\textbf{18.} Lots espérés par billet~:
$\frac{100 + 2 \times 50}{200} = 1$ euro contre un billet à $2$
euros~: marge de $50\,\%$ --- dix-huit fois les $2.7\,\%$ de la
roulette. La tombola s'en tire parce que ses joueurs font sciemment
un don~: la «~perte~» est précisément le but.

\textbf{19.} Sinistres espérés~: $1\,000 \times 0.01 \times
10\,000 = 100\,000$~; primes~: $130\,000$~: bénéfice espéré
$30\,000$ euros. \omterm{def:g11:stat:variance}{Écart type} du total des sinistres~:
$10\,000 \times \sqrt{1\,000 \times 0.01 \times 0.99}
\approx 31\,500$ euros --- une seule mauvaise année ordinaire
dévore tout le bénéfice espéré. D'où les réserves de capital, la
réassurance et des portefeuilles bien plus vastes que mille
contrats~: les assureurs vivent de la loi des grands nombres et
gardent des réserves contre sa lenteur.

\textbf{20.} La liste de contrôle d'abord --- $n$ fixé, même $p$,
\omterm{def:g12:condprob:indep}{indépendance} \emph{revendiquée} --- sans quoi il n'y a pas de
binomiale du tout. Ensuite, naviguer avec $\E \pm \sigma$~: les
\omterm{def:g11:stat:mean}{moyennes} situent, les écarts avertissent. Puis tarifer les queues~:
refus d'embarquement, mauvais lots et années ruineuses vivent tous
au-delà de $2\sigma$. Réserves~: l'\omterm{def:g12:condprob:indep}{indépendance} est une promesse du
modélisateur, pas du monde~; et l'\omterm{def:g12:randvar:exp}{espérance} ignore précisément ce
qui vous détruit. Les pages manquantes du règlement --- à quelle
vitesse les \omterm{def:g10:stats:series}{fréquences} se stabilisent, et quelle forme prennent les
fluctuations --- sont les deux chapitres suivants.
\end{solution}
